% from ega4-11-fr.tex
\oldpage[IV-3]{146}
Pour qu'un homomorphisme $u:A\to A'$ soit tel que $M\otimes_A A'$ soit un
$A'$-module libre \textup{(resp.} un $A'$-module plat), il faut et il suffit que $u$ se
factorise en $A\to A/\mathfrak K_0\to A'$ \textup{(ou encore que
$\mathfrak K_0A'=0$)}.
\end{corollary}

\begin{proof}
Le fait que l'intersection $\mathfrak K_0$ des idéaux $\mathfrak K$ pour
lesquels $M\otimes_A(A/\mathfrak K)$ est un $(A/\mathfrak K)$-module libre
soit le plus petit de ces idéaux résulte de (11.4.1) appliqué à l'anneau
$A/\mathfrak K_0$, à son idéal nilpotent $\mathfrak J/\mathfrak K_0$ et aux
homomorphismes $A/\mathfrak K_0\to A/\mathfrak K$, dont les noyaux ont $0$
pour intersection. Si $A'$ est une $(A/\mathfrak K_0)$-algèbre, on a
$M\otimes_A A'=
(M\otimes_A(A/\mathfrak K_0))\otimes_{A/\mathfrak K_0}A'$,
donc $M\otimes_A A'$ est un $A'$-module libre. Réciproquement, si $A'$ est
une $A$-algèbre telle que $M\otimes_A A'$ soit un $A'$-module libre et si
$\mathfrak K$ est le noyau de l'homomorphisme $A\to A'$, il résulte de
(11.4.1) appliqué à l'anneau $A/\mathfrak K$, au
$(A/\mathfrak K)$-module $M\otimes_A(A/\mathfrak K)$, à l'idéal nilpotent
$\mathfrak J\mathfrak K/\mathfrak K$ et à l'homomorphisme \emph{injectif}
$A/\mathfrak K\to A'$, que $M\otimes_A(A/\mathfrak K)$ est un
$(A/\mathfrak K)$-module libre, donc que $\mathfrak K\supset\mathfrak K_0$.
Le fait que l'on puisse remplacer «~libre~» par «~plat~» dans ce qui précède
(en conservant naturellement l'hypothèse que $M\otimes_A(A/\mathfrak J)$ est
un $(A/\mathfrak J)$-module \emph{libre}) résulte de
$(\mathbf 0_{\mathrm{III}},10.1.2)$, comme on l'a vu au début de la
démonstration de (11.4.1).
\end{proof}

\begin{proposition}[11.4.5]
\label{IV.11.4.5-fr}
Soient $Y$ un préschéma irréductible, $f:X\to Y$ un morphisme de présentation
finie, $\mathcal F$ un $\mathcal O_X$-Module de présentation finie. Il existe
alors un ouvert non vide $U$ dans $Y$ et un sous-schéma fermé $Z$ de $U$, de
présentation finie sur $U$, ayant la propriété suivante~: pour tout changement
de base $U'\to U$, si l'on pose $X'=X\times_YU'$, $f'=f_{(U')}$ et
$\mathcal F'=\mathcal F\otimes_{\mathcal O_Y}\mathcal O_{U'}$, alors, pour que
$\mathcal F'$ soit $f'$-plat, il faut et il suffit que le morphisme $U'\to U$
se factorise en $U'\to Z\to U$. Un tel schéma $Z$ a même espace sous-jacent que
$U$. Supposons de plus $Y$ affine, et soit $(W_i)$ un recouvrement fini de $X$
par des ouverts affines~; alors, on peut supposer $U$ choisi de sorte que, si
$U'=\operatorname{Spec}(A')$ est un schéma affine et $U'\to U$ un morphisme qui
se factorise en $U'\to Z\to U$, chacun des
$\Gamma(W_i\times_YU',\mathcal F')$ est un $A'$-module libre.
\end{proposition}

\begin{proof}
On peut évidemment se borner au cas où $Y=\operatorname{Spec}(A)$ est affine.
Utilisant (8.9.1), il y a un sous-anneau noethérien $A_1$ de $A$, un morphisme
de type fini $f_1:X_1\to Y_1=\operatorname{Spec}(A_1)$ et un
$\mathcal O_{X_1}$-Module cohérent $\mathcal F_1$ tels que $f=f_{(Y)}$ et
$\mathcal F=\mathcal F_1\otimes_{\mathcal O_{Y_1}}\mathcal O_Y$~; on peut en
outre supposer que les $W_i$ sont images réciproques d'ouverts affines de
$X_1$. Notons en outre que $Y_1$ est irréductible, le morphisme $Y\to Y_1$
étant dominant $(\mathbf I,1.2.7)$. Supposons la proposition démontrée pour
$Y_1$, $X_1$ et $\mathcal F_1$, et soient $U_1$ l'ouvert de $Y_1$ et $Z_1$ le
sous-schéma fermé de $U_1$ ayant les propriétés voulues, $U=U_1\times_{Y_1}Y$
et $Z=Z_1\times_{Y_1}Y$ leurs images réciproques. Alors, si $U'\to U$ est un
changement de base tel que
$\mathcal F'=\mathcal F\otimes_{\mathcal O_Y}\mathcal O_{U'}
=\mathcal F_1\otimes_{\mathcal O_{Y_1}}\mathcal O_{U'}$ soit $f'$-plat, le
morphisme $U'\to U_1$ se factorise en $U'\to Z_1\to U_1$~; mais comme
$U'\to U_1$ se factorise aussi en $U'\to U\to U_1$, la définition du produit
de préschémas montre que $U'\to U$ se factorise en $U'\to Z\to U$.

On peut donc se borner au cas où $A$ est \emph{noethérien}~; soit
$\mathfrak N$ son nilradical, qui est nilpotent, et posons
$A_0=A_{\mathrm{red}}=A/\mathfrak N$,
$Y_0=\operatorname{Spec}(A_0)=Y_{\mathrm{red}}$,
$X_0=X\times_YY_0$, $\mathcal F_0=\mathcal F\otimes_{\mathcal O_Y}
\mathcal O_{Y_0}$, $W_{i0}=W_i\times_YY_0$~; si
$M_i=\Gamma(W_i,\mathcal F)$, $M_{i0}=\Gamma(W_{i0},\mathcal F_0)$ est donc
égal à $M_i\otimes_A A_0$. Comme $A_0$ est \emph{intègre}, on peut, en vertu du
théorème de platitude générique (6.9.2), en remplaçant au besoin $Y$ par un
ouvert non vide de $Y$, supposer que les $M_{i0}$ soient des $A_0$-modules
\emph{libres}. En vertu de (11.4.4), il y a donc pour chaque $i$ un idéal
$\mathfrak J_i\subset\mathfrak N$ tel
